% Part: normal-modal-logic
% Chapter: tableaux
% Section: simple-S5

\documentclass[../../../include/open-logic-section]{subfiles}

\begin{document}

\olfileid{nml}{tab}{s5}

\olsection{\usetoken{P}{tableau} مبسطة للنسق \Log{S5}}

للنسق~\Log{S5} السلامة والتمام بالنسبة إلى النماذج الكلية، وهي التي
ينفذ فيها كل عالم إلى كل عالم. فلا يبقى لتعيين علاقة الإتاحة أثر زائد:
وجود عالم~${\OLId{latin-ordinary}{w}}$ يحقق~$\mSat{{\OLId{latin-looped}{M}}}{!A}[{\OLId{latin-ordinary}{w}}]$ يكافئ وجود عالم كذلك يمكن
النفاذ إليه من~${\OLId{latin-ordinary}{u}}$. فيجوز، عند دراسة~\Log{S5}، أن نكتفي في وصف
النموذج بمجموعة عوالم وتقويم~${\OLId{latin-looped}{V}}$. ويتبع ذلك تبسيط قواعد ال!!{tableau}.
فإنما احتجنا، في الحالة العامة، إلى متتاليات الأعداد الصحيحة الموجبة
كي نستدل من البادئات على النفاذ؛ فـ~$\sigma.{\OLId{latin-ordinary}{n}}$ تسمّي عالمًا ينفذ إليه
العالم الذي تسميه~$\sigma$. أما في~\Log{S5}، بالنماذج الكلية، فالنفاذ
عام بين العوالم، فلا حاجة إلى حفظ هذه الصلات في أسماء البادئات.
ونستعيض عن المتتاليات بأعداد صحيحة موجبة مفردة. والقواعد على هذا الوجه
في~\olref{tab:rules-S5}.

\begin{table}
  \begin{center}
    \def\arraystretch{3}\def\fCenter{}
    \begin{tabular}{|c|c|}
    \hline
    \iftag{prvBox}{
      \AxiomC{\sFmla{\True}{\Box !A}[n]}
      \RightLabel{\TRule{\True}{\Box}}
      \UnaryInfC{\sFmla{\True}{!A}[m]}
      \DisplayProof
      &
      \AxiomC{\sFmla{\False}{\Box !A}[n]}
      \RightLabel{\TRule{\False}{\Box}}
      \UnaryInfC{\sFmla{\False}{!A}[m]}
      \DisplayProof\\
      ${\OLId{latin-ordinary}{m}}$ مستعملة & ${\OLId{latin-ordinary}{m}}$ جديدة
      \\[1ex]
      \hline}{}
    \iftag{prvDiamond}{
      \AxiomC{\sFmla{\True}{\Diamond !A}[n]}
      \RightLabel{\TRule{\True}{\Diamond}}
      \UnaryInfC{\sFmla{\True}{!A}[m]}
      \DisplayProof
      &
      \AxiomC{\sFmla{\False}{\Diamond !A}[n]}
      \RightLabel{\TRule{\False}{\Diamond}}
      \UnaryInfC{\sFmla{\False}{!A}[m]}
      \DisplayProof\\
      ${\OLId{latin-ordinary}{m}}$ جديدة & ${\OLId{latin-ordinary}{m}}$ مستعملة
      \\[1ex]
      \hline}{}
    \end{tabular}
  \end{center}
  \caption{قواعد مبسطة للنسق \Log{S5}.}
  \ollabel{tab:rules-S5}
\end{table}

\begin{ex}
  لإثبات~$\Log{S5} \Proves
  \Ax{5}$، أي~$\Diamond!A \lif \Box\Diamond!A$، يكفينا الجدول الدلالي
  المبسط الآتي، وهو مغلق.
  \begin{oltableau}
    [\pFmla{\False}{\Diamond\formula{A} \lif \Box\Diamond \formula{A}}{1},
      just = \TAss
      [\pFmla{\True}{\Diamond \formula{A}}{1}, just = {\TRule{\False}{\lif}[1]}
        [\pFmla{\False}{\Box\Diamond \formula{A}}{1},
          just = {\TRule{\False}{\lif}[1]}
          [\pFmla{\False}{\Diamond \formula{A}}{2},
            just = {\TRule{\False}{\Box}[3]}
            [\pFmla{\True}{\formula{A}}{3},
              just = {\TRule{\True}{\Diamond}[2]}
                [\pFmla{\False}{\formula{A}}{3},
                  just = {\TRule{\False}{\Diamond}[4]}, close]
            ]
          ]
        ]
      ]
    ]
  \end{oltableau}
\end{ex}

\end{document}
